\chapter{Entropy maximizers and translation symmetries}\label{sec:shrinker-entropy}
A self-shrinker's normalized Gaussian area is its entropy. To prove this on a noncompact surface, we keep the Gaussian integration and equality cases explicit. Our reference is \cite{CM12}, Lemma 7.10 and equations (7.14)--(7.27), pp.~785--788. The shrinker normalization is $\vec H=-x^\perp/2$. We write $[u]_{c,t}=\int_\Sigma u\rho_{c,t}\,dA$.

\begin{lemma}[Shifted weighted divergence identities]\label{lem:shrinker-shifted-divergence}
Let $\Sigma^n\subset\mathbb R^{n+1}$ be a smooth complete properly immersed boundaryless self-shrinker with polynomial area growth. For any $c,y\in\mathbb R^{n+1}$ and $t>0$,
\[
 [\langle(x-c)^T,y^T\rangle]_{c,t}
 =-t[\langle x^\perp,y^\perp\rangle]_{c,t},
\]
\[
 [\langle(x-c)^T,x^T\rangle]_{c,t}
 =t[2n-|x^\perp|^2]_{c,t}.
\]
\end{lemma}
\begin{proof}
The shifted drift $\mathcal D_{c,t}=\Delta-\langle(x-c)^T,\nabla\cdot\rangle/(2t)$ is $\rho_{c,t}^{-1}\operatorname{div}_\Sigma(\rho_{c,t}\nabla\cdot)$. Since $\Delta x=-x^\perp/2$,
\[
 2\mathcal D_{c,t}\langle x,y\rangle
 =-\langle x^\perp,y^\perp\rangle-\frac{\langle(x-c)^T,y^T\rangle}{t},
\]
\[
 \mathcal D_{c,t}|x|^2
 =2n-|x^\perp|^2-\frac{\langle(x-c)^T,x^T\rangle}{t}.
\]
The integral of each drift derivative is zero. Indeed, test its weighted divergence with compact radial cutoffs. The error is bounded by a polynomial times a Gaussian on the receding annulus, and tends to zero by Lemma \ref{lem:gaussian-polynomial-moments}; properness makes these cutoffs compact. Integrating the displayed identities proves the claims. This is \cite{CM12}, (7.18)--(7.19).
\end{proof}

\begin{proposition}[Gaussian area decreases along parabolic parameter paths]\label{prop:shrinker-parameter-monotonicity}
Let $\Sigma$ satisfy the hypotheses of Lemma \ref{lem:shrinker-shifted-divergence}. Fix $y\in\mathbb R^{n+1}$ and $a\in\mathbb R$, set $t_s=1+as^2$, and put $g(s)=F_{sy,t_s}(\Sigma)$ whenever $t_s>0$. Then for $s\geq0$,
\[
 g'(s)=-\frac{s}{2t_s^2}
 \bigl[|(asx+y)^\perp|^2\bigr]_{sy,t_s}\leq0.
\]
\end{proposition}
\begin{proof}
Write $z=x-sy$ and suppress the weighted subscript. The preceding identities imply
\[
 \left[\frac{\langle z^T,y^T\rangle}{t_s}\right]
 =-[\langle z^\perp,y^\perp\rangle+s|y^\perp|^2],
\]
\[
 \left[\frac{|z^T|^2}{t_s}\right]
 =[2n-|z^\perp|^2-s\langle z^\perp,y^\perp\rangle].
\]
For the second equation substitute $x=z+sy$ into the second shifted identity and use the first. Differentiation under the Gaussian integral is justified by polynomial growth on any compact interval with $t_s>0$. It gives
\[
 g' =\frac{t_s'}{4t_s}
 \left[\frac{|z^T|^2}{t_s}+\frac{|z^\perp|^2}{t_s}-2n\right]
 +\frac12\left[\frac{\langle z^T,y^T\rangle}{t_s}
 +\frac{\langle z^\perp,y^\perp\rangle}{t_s}\right].
\]
Insert the two identities and use $t_s'=2as$ and $t_s^{-1}-1=-as^2/t_s$. Completing the square yields
$g'=-s[|as z^\perp+t_s y^\perp|^2]/(2t_s^2)$.
Since $as z+t_s y=asx+y$, this is the stated formula. It is \cite{CM12}, (7.27), with the square simplified.
\end{proof}

\begin{theorem}[A shrinker attains its entropy at its normalized parameters]\label{thm:shrinker-entropy-maximum}
For a smooth complete properly immersed boundaryless self-shrinker $\Sigma^n\subset\mathbb R^{n+1}$ with polynomial area growth,
\[
 F_{c,t}(\Sigma)\leq F_{0,1}(\Sigma),\qquad
 \lambda(\Sigma)=F_{0,1}(\Sigma)<\infty
\]
for all $c$ and $t>0$. Consequently
\[
 \operatorname{Area}(\Sigma\cap B_r(c))
 \leq e^{1/4}(4\pi)^{n/2}F_{0,1}(\Sigma)r^n.
\]
This is the maximum statement proved in \cite{CM12}, Section 7.2.
\end{theorem}
\begin{proof}
In Proposition \ref{prop:shrinker-parameter-monotonicity}, choose $y=c$ and $a=t-1$. Along $0\leq s\leq1$, $t_s=(1-s^2)+s^2t$ is positive. Hence $F_{c,t}=g(1)\leq g(0)=F_{0,1}$. The latter is finite by Lemma \ref{lem:gaussian-polynomial-moments}. Take the supremum to get equality with entropy. Proposition \ref{prop:mcf-entropy-area-growth} gives the uniform ball-area bound. Thus the initially assumed polynomial exponent need not equal $n$; the shrinker equation improves it to the displayed bound.
\end{proof}

\begin{proposition}[Constant tangent vectors and product splitting]\label{prop:shrinker-translation-splitting}
For a connected complete properly embedded smooth hypersurface $\Sigma\subset\mathbb R^{n+1}$, define
\[
 E=\{v\in\mathbb R^{n+1}:v\in T_x\Sigma\text{ for every }x\in\Sigma\}.
\]
This is a linear subspace. The surface is invariant under translation by $E$ and splits orthogonally as
\[
 \Sigma=(\Sigma\cap E^\perp)\times E.
\]
In particular, a nonzero constant vector is tangent everywhere exactly when the surface splits off the corresponding Euclidean line.
\end{proposition}
\begin{proof}
The intersection of the linear spaces $T_x\Sigma$ is linear. A constant tangent vector $v$ induces a smooth bounded vector field on the complete induced manifold. Its integral curves exist for all real time: on every bounded time interval their lengths are bounded, so completeness and local ODE existence prevent a finite endpoint. Their ambient derivatives equal $v$, hence they are $x+sv$. This proves translation invariance in each direction of $E$.

Translate $x\in\Sigma$ by $-\pi_Ex$ to obtain $x-\pi_Ex\in\Sigma\cap E^\perp$. The converse inclusion follows by translation invariance. The restriction of $\pi_E$ to $\Sigma$ is a submersion because $E\subset T_x\Sigma$; its zero fiber is a smooth submanifold. The map from this fiber times $E$ to $\Sigma$ by addition is a smooth bijection with smooth inverse $(\pi_{E^\perp},\pi_E)$, and orthogonality makes the induced metric a product. No shrinker equation is needed for this lemma.
\end{proof}

\begin{lemma}[A smooth proper minimal shrinker is a plane]\label{lem:minimal-shrinker-plane}
If a connected complete properly embedded smooth boundaryless self-shrinker has $\vec H=0$ everywhere, it is a linear hyperplane through the origin.
\end{lemma}
\begin{proof}
The shrinker equation gives $x^\perp=0$, so the position field is tangent. Its integral curve through $p$ is $e^sp$. It exists for every finite positive or negative time: its speed on any finite interval is bounded, and completeness supplies the continuation argument used in Proposition \ref{prop:shrinker-translation-splitting}. Thus $rp\in\Sigma$ for all $r>0$. Properness makes the embedded image closed; letting $r\downarrow0$ gives $0\in\Sigma$.

Smoothness at zero implies that normalized secants approaching zero approach its tangent plane. Applying this to $rp$ shows $p\in T_0\Sigma$ for every $p\in\Sigma$. The inclusion of $\Sigma$ in the $n$-plane $T_0\Sigma$ is locally open by full rank. It is closed by properness and nonempty, so connectedness of the plane makes it the entire plane. This supplies the smooth-cone step used in \cite{CM12}, p.~786.
\end{proof}

\begin{proposition}[All maximizing parameters of an embedded shrinker]\label{prop:shrinker-gaussian-equality}
Let $\Sigma$ be a connected complete smooth properly embedded boundaryless self-shrinker with polynomial area growth, and let $E$ be the constant tangent subspace of Proposition \ref{prop:shrinker-translation-splitting}. If $\Sigma$ is not a plane, then
\[
 F_{c,t}(\Sigma)=\lambda(\Sigma)\quad\Longleftrightarrow\quad t=1\text{ and }c\in E.
\]
For a linear hyperplane $P$, equality holds precisely when $c\in P$, with any $t>0$.
\end{proposition}
\begin{proof}
Use the path in Theorem \ref{thm:shrinker-entropy-maximum}. Equality at both endpoints forces its nonincreasing smooth function to be constant. For every $0<s<1$, the integral square in Proposition \ref{prop:shrinker-parameter-monotonicity} is zero. The weight is positive and the integrand continuous, so
$((t-1)sx+c)^\perp=0$ at every point. Subtract this identity at two different $s$. If $t\ne1$, it gives $x^\perp=0$, and Lemma \ref{lem:minimal-shrinker-plane} forces a plane. Otherwise $c^\perp=0$, which is $c\in E$. Conversely, if $t=1$ and $c\in E$, the square vanishes all along the path and gives equality. The plane statement follows from Example \ref{ex:mcf-plane-entropy}. This identifies the degeneracy excluded in \cite{CM12}, Lemma 7.10.
\end{proof}

\begin{proposition}[A quadratic parameter deficit without line splitting]\label{prop:shrinker-local-entropy-gap}
Let $\Sigma$ be a connected complete smooth properly embedded two-sided boundaryless self-shrinker with polynomial area growth and with $E=\{0\}$. There are $\gamma,r>0$ such that
\[
 F_{c,t}(\Sigma)\leq\lambda(\Sigma)-\gamma\bigl(|c|^2+(t-1)^2\bigr)
\]
whenever $|c|^2+(t-1)^2\leq r^2$ and $t>0$.
\end{proposition}
\begin{proof}
The first parameter derivatives vanish by Proposition \ref{prop:gaussian-joint-first}. With the shape fixed, Theorem \ref{thm:shrinker-joint-hessian} gives the Hessian on $(y,h)$ as
\[
 -h^2\int H^2\rho_{0,1}\,dA-\tfrac12\int|y^\perp|^2\rho_{0,1}\,dA.
\]
Since $E=0$, every nonzero $y$ gives a positive second integral. Also $H$ is not identically zero: Lemma \ref{lem:minimal-shrinker-plane} would give a plane, whose tangent subspace has dimension $n\geq1$. The Hessian is therefore negative definite. Compactness of the unit sphere in the finite-dimensional parameter space gives a uniform negative bound. Gaussian domination gives continuity of second derivatives near $(0,1)$, so Taylor's theorem proves the quadratic estimate after decreasing $r$. Compare \cite{CM12}, (7.12).
\end{proof}

\begin{theorem}[A uniform entropy gap away from the normalized parameters]\label{thm:shrinker-global-entropy-gap}
Let $\Sigma$ satisfy the hypotheses of Proposition \ref{prop:shrinker-local-entropy-gap}. For every $\varepsilon>0$ there is $\delta>0$ such that
\[
 \sup_{|c|+|\log t|\geq\varepsilon}F_{c,t}(\Sigma)
 \leq\lambda(\Sigma)-\delta.
\]
This is the uniform, rather than merely pointwise, conclusion of \cite{CM12}, Lemma 7.10. In particular $\lambda(\Sigma)>1$.
\end{theorem}
\begin{proof}
Choose a closed Euclidean parameter ball of radius $r_0>0$ about $(0,1)$, small enough to lie in the quadratic-deficit neighborhood and in $\{|c|+|\log t|<\varepsilon\}$. Its boundary has deficit at least $\gamma r_0^2$. For any target outside $\{|c|+|\log t|<\varepsilon\}$, the continuous path $(sc,1+(t-1)s^2)$ starts at the center and ends outside this ball. It crosses the boundary at some $s_0\in(0,1]$. Its scale remains positive, and Proposition \ref{prop:shrinker-parameter-monotonicity} makes the Gaussian area nonincreasing after that crossing. Thus every target has deficit at least $\delta=\gamma r_0^2$.

Choose a point $p\in\Sigma$ and let $t\downarrow0$. The parameters $(p,t)$ eventually satisfy the excluded-neighborhood condition. Proposition \ref{prop:static-gaussian-density} gives $F_{p,t}\to1$, so $1\leq\lambda(\Sigma)-\delta$. This proves strict entropy greater than one without a separate classification theorem.
\end{proof}
